% Generated draft: requires mathematical and source review.
\documentclass[11pt]{article}
\usepackage[margin=25mm]{geometry}
\usepackage{fontspec}
\setmainfont{Latin Modern Roman}
\usepackage{amsmath,amssymb,mathtools,unicode-math}
\setmathfont{Latin Modern Math}
\usepackage{enumitem,needspace,xurl,hyperref}
\setlength{\parindent}{0pt}
\setlength{\parskip}{5pt}
\setlength{\emergencystretch}{3em}
\newcommand{\N}{\mathbb{N}}
\newcommand{\Z}{\mathbb{Z}}
\newcommand{\Q}{\mathbb{Q}}
\newcommand{\R}{\mathbb{R}}
\newcommand{\C}{\mathbb{C}}
\newcommand{\F}{\mathbb{F}}
\newcommand{\A}{\mathbb{A}}
\newcommand{\PP}{\mathbb{P}}
\newcommand{\E}{\mathbb{E}}
\newcommand{\eps}{\varepsilon}
\newcommand{\dd}{\,\mathrm{d}}
\begin{document}
\section*{UnsolvedMath ID: 20002333}
\textbf{Source problem code:} catalogue entry
\subsection{Definitions and mathematical statement}
The original problem is to study definable sets in an expansion of Urysohn space by an involution. The extract does not specify bounded versus unbounded Urysohn space, the choice of involution, allowed parameters or sorts, or whether ``definable set'' means a zero set or a set with definable distance. These choices affect the question.

Here is the particular setting of the supplied research aid. Let $U$ be the complete separable Urysohn sphere of diameter $1$, and let $\sigma:U\to U$ be an isometric involution, $\sigma^2=\operatorname{id}$. Naming this map by a function symbol is equivalent to naming its metric graph predicate $P(x,y)=d(\sigma x,y)$. Assume the exact two-extension property: whenever $A\subseteq U$ is finite and $\sigma$-invariant, every diameter-at-most-one finite metric extension of $A$ by one involution orbit of size one or two, with an isometric involution extending $\sigma|_A$, embeds equivariantly into $U$ over $A$. This specifies the generic involution setting used below; no assertion is made about an arbitrary named involution.

A unary definable predicate without parameters is a uniform limit of formula predicates without parameters. A nonempty closed subset $D\subseteq U$ is distance-definable without parameters when $x\mapsto d(x,D)=\inf_{y\in D}d(x,y)$ is such a predicate. Write $\delta(x)=d(x,\sigma x)$. The aid reports the following classification in the home sort:
\[
 \{\hbox{unary definable predicates without parameters}\}
 =\{f\circ\delta:f:[0,1]\to[0,1]\text{ continuous}\},
\]
and the nonempty distance-definable sets without parameters are exactly
\[
 D_K=\{x\in U:\delta(x)\in K\},
 \qquad\varnothing\ne K\subseteq[0,1]\text{ closed}.
\]
The exact claimed distance formula is
\[
 d(x,D_K)=\tfrac12\operatorname{dist}(\delta(x),K),
 \qquad \operatorname{dist}(t,K)=\inf_{s\in K}|t-s|.
\]
Each level $D_{\{r\}}$, $0\le r\le1$, is reported to be isometric to the Urysohn sphere. However, for every $0<r\le1$, the distance-definable set $D_{\{0,r\}}$ is reported not to be Urysohn: it contains $a,b$ with $\delta(a)=0$, $\delta(b)=r$ and $d(a,b)=r/2$, having no midpoint in that set. Indeed $\delta$ is $2$-Lipschitz; a midpoint with displacement $r$ is excluded by its distance $r/4$ from $a$, while a midpoint with displacement $0$ is excluded by its distance $r/4$ from $b$.

These statements address unary predicates and sets without parameters for this generic isometric involution. They do not classify sets in powers, imaginary sorts, expansions with parameters, or other involutions. The aid describes a counterexample to a historical conjecture under this restricted reading, but the conjecture itself is not included in the supplied original extract.
\subsection{Short English statement}
Study definable sets after naming an involution of Urysohn space. In the research aid's bounded generic setting, unary sets without parameters are classified by closed displacement conditions; individual levels are Urysohn, while a union of the zero level with a positive level need not be.
\subsection{Sources}
\begin{itemize}
\item American Institute of Mathematics, workshop \emph{Model theory of metric structures}, Problem 19, \url{https://aimath.org/WWN/continuouslogic/continuouslogic.pdf}.
\item Local UnsolvedMath registry, permanent ID 20002333, source code \texttt{catalogue entry}, canonical location \texttt{aim-logic-notes.json}, \texttt{notes[108]}, record dated 14 August 2026. The extracted ``Background:'' contains no further text, and the original literature field is empty.
\item Supplied research attempt 1: counterexample, problem status at run partially solved, rigor labelled proved, novelty confidence low. Its displacement classification and two-level obstruction are machine-generated research aid, not peer reviewed; independently verify before citation. No separate authoritative reference for these refinements is supplied.
\end{itemize}
\end{document}
